% Generated by tools/paper/render.py from main_text.py; edit the source there, not this file.
\section{Materials and methods}

\subsection{Overview}

pgphase takes the reads' alignments to a pangenome graph, as a GAF file \citep{Li2020}, the graph's site catalog, and optionally the same reads aligned to the linear reference, as a BAM file (Fig.~\ref{fig:overview}). It works on chunks of the reference (1~Mb by default with a BAM), in parallel, and joins the chunks through the reads they share. Its two parts are the \emph{graph arm}, which phases the catalog's sites from the graph alignments, and \emph{union gap phasing}, which adds the sample's own heterozygous variants from the linear alignments and phases everything in one solve. Its outputs are a phased VCF and the reads' haplotype (HP) and phase set (PS) tags.

% Generated by tools/paper/render.py from main_text.py; edit the source there, not this file.
\begin{figure*}[t]
\centering
\resizebox{\textwidth}{!}{\begin{tikzpicture}
\node[inbox, minimum width=2.6cm, minimum height=0.9cm, anchor=west] (graph) at (0.0,4.2) {Pangenome graph\\(GBZ)};
\node[inbox, minimum width=2.6cm, minimum height=0.9cm, anchor=west] (reads) at (0.0,2.3) {Long reads};
\node[box, minimum width=2.8cm, minimum height=0.9cm, anchor=west] (catalog) at (3.4,4.2) {Site catalog\\(snarl alleles, pgcat)};
\node[box, minimum width=2.8cm, minimum height=0.9cm, anchor=west] (gaf) at (3.4,2.9) {Graph alignments\\(GAF, indexed)};
\node[box, minimum width=2.8cm, minimum height=0.9cm, anchor=west] (bam) at (3.4,1.6) {Linear alignments\\(BAM)};
\node[box, minimum width=3.2cm, minimum height=1.2cm, anchor=west] (clean) at (7.0,3.55) {Catalog sites phased\\from walk observations\\(k-means)};
\node[box, minimum width=3.2cm, minimum height=1.2cm, anchor=west] (solve) at (7.0,1.6) {Sample-specific hets\\(candidates, noisy-region\\MSA)};
\node[box, minimum width=3.3cm, minimum height=1.6cm, anchor=west] (em) at (11.0,2.6) {Global EM per chunk:\\site phases and errors,\\joins the reads verify};
\node[box, minimum width=3.0cm, minimum height=1.0cm, anchor=west] (windows) at (15.1,3.45) {Haplotype and allele\\windows, EM pass 2};
\node[box, minimum width=3.0cm, minimum height=1.0cm, anchor=west] (fill) at (15.1,1.75) {Stitching, rescue,\\island blocks};
\node[outbox, minimum width=2.4cm, minimum height=1.3cm, anchor=west] (out) at (18.9,2.6) {Phased VCF\\Haplotagged\\reads (HP/PS)};
\draw[arr] (2.60,4.20) -- (3.40,4.20);
\draw[arr] (2.60,2.30) -- (3.40,2.90);
\draw[arr] (2.60,2.30) -- (3.40,1.60);
\draw[arr] (2.60,4.20) -- (3.40,2.90);
\draw[arr] (6.20,4.20) -- (7.00,3.55);
\draw[arr] (6.20,2.90) -- (7.00,3.55);
\draw[arr] (6.20,1.60) -- (7.00,1.60);
\draw[arr] (10.20,3.55) -- (11.00,2.60);
\draw[arr] (10.20,1.60) -- (11.00,2.60);
\draw[arr] (14.30,2.60) -- (15.10,3.45);
\draw[arr] (16.60,2.95) -- (16.60,2.25);
\draw[arr] (14.30,2.60) -- (15.10,1.75);
\draw[arr] (18.10,1.75) -- (18.90,2.60);
\node[font=\scriptsize\itshape, anchor=west] at (7.0,5.0) {graph arm: known alleles};
\node[font=\scriptsize\itshape, anchor=west] at (7.0,0.7) {union gap phasing: the sample's own hets};
\end{tikzpicture}}
\caption{pgphase's pipeline. The graph arm phases the sites of a pangenome graph's catalog from the walks the reads' graph alignments take through them; with a linear alignment, union gap phasing adds the sample's own heterozygous variants, and one expectation-maximization phases all sites, joining blocks only where reads verify the join. Reads still unlabelled are phased from local windows, rescue and island blocks.}\label{fig:overview}
\end{figure*}


\subsection{Site catalog and read evidence from a pangenome graph}

The catalog is built once per graph (\texttt{build-snarl-catalog}, from a GBZ graph \citep{Siren2022}): each snarl \citep{Paten2018} of the graph is a site, its reference allele the walk of the reference path through it and its alternative alleles the distinct walks of the haplotype paths, ordered by how many haplotypes take them. The reads' graph alignments are indexed once by reference position (\texttt{index-gaf}). For each chunk, every read crossing a site contributes the allele whose walk its own alignment follows; the identity of an allele is its walk, so two reads carry the same allele exactly when they take the same path through the site. A multi-allelic site is split into one reference-against-alternative contrast per alternative allele, and sites are kept by depth and allele fraction. Indels in homopolymers, tandem repeats and low-complexity sequence are set apart as noisy candidates, which do not take part in the first solve.

\subsection{Phasing the catalog sites}

The catalog's clean heterozygous sites are phased first by k-means clustering of the reads on their alleles, the algorithm of longcallD \citep{Gao2026} applied to walk observations. With a BAM, each graph SNP is checked against the bases of well-mapped reads: one whose reference base no such read shows, and no deletion explains, is reclassified homozygous. Where a read takes a walk that differs from a selected allele only away from a SNP, the walk is read as that SNP's allele, so that reads whose alignment takes a rarer path through a site still vote. The clean solve's labels start the joint solve below.

\subsection{The sample's own heterozygous variants}

With a BAM, each chunk is also solved from its linear alignments, by the alignment pipeline pgphase ports from longcallD \citep{Gao2026}: candidate variants from the reads' differences to the reference, classified by depth, allele fraction and sequence context, with noisy regions up to 2~kb resolved by aligning each haplotype's reads with partial order alignment \citep{Gao2021}. Its clean heterozygous variants, and those verified by the multiple sequence alignment, join the chunk's sites with their per-read alleles. A verified noisy variant must be seen by two reads of mapping quality 30 on each of its two alleles, and no injected site may be read by more than 1.5 times the chunk's median depth: at a collapsed duplication the two "alleles" are the copies, and a run of such sites agrees with itself perfectly. Variants that fail these rules still label reads, but only as islands (Section 2.7). The same indel placed at several copies of a repeat becomes one site; calls inside a catalog site of up to 2~kb only give alleles to reads the graph left without one there; reads only the linear alignment holds are added. At injected indels, every read crossing the site without a call is realigned to the two alleles in reference context and takes the closer one, with a per-read error of \(1/(1+e^{3m})\) for a margin of \(m\) edits. At catalog SNPs a read's graph allele is replaced by its aligned base where the two disagree: measured against the parents, the alignment's base was right in about 90\% of such conflicts.

\subsection{One expectation-maximization over all sites}

Every site with two observed alleles, from the catalog or the sample, gets a phase \(h_s\in\{0,1\}\) and an error rate \(\varepsilon_s\), and every read \(r\) a posterior probability of each haplotype. With \(x_{rs}\) the allele read \(r\) shows at site \(s\) and \(e_{rs}\) the per-read error of that call where one is known, a read on haplotype \(a\) has likelihood

\begin{equation}
L_r(a) = \prod_{s} \bigl(1-\varepsilon_{rs}\bigr)^{[x_{rs} = h_s \oplus a]}\, \varepsilon_{rs}^{[x_{rs} \neq h_s \oplus a]},
\end{equation}

where \(\varepsilon_{rs} = 1-(1-\varepsilon_s)(1-e_{rs})\) combines the site's error with the call's, and expectation-maximization \citep{Dempster1977} alternates the reads' posteriors with the sites' phases and errors. A site's error is the weighted disagreement of its reads, shrunk toward the median error of the sites of its class in the chunk (clean catalog heterozygotes, other catalog sites, injected sites and local windows, each split by whether the site is a one-base length difference) by five pseudo-reads, an empirical Bayes prior that keeps a noisy repeat call from overruling a clean SNP that shares its reads. Only reads of mapping quality 20 or more teach the model; the others are scored against it and labelled. A site whose rarer allele is held by fewer than a quarter of its reads is left out: a skewed split is the signature of a paralog or an error. After the first iterations, the solve tries switch moves, flipping every site past the boundary whose flip most increases the likelihood. A site whose learned error is 0.3 or more is reported unphased and labels no read; a read is labelled in the block where its own evidence is strongest when its posterior there is at least 0.8.

\subsection{Joining blocks only where reads verify the join}

A block ends wherever flipping everything after a boundary would cost less than 4 log-likelihood units, or fewer than two reads carry evidence on both sides: one molecule can be chimeric or misplaced. Beyond that, pgphase restricts which sites can carry a join at all, because a wrong join flips every read past it:

\begin{enumerate}
  \item \textbf{One-base length differences are weak.} A homopolymer length call in a homozygous stretch looks like a heterozygous site; a join that holds only through such sites needs a margin of 10 log units instead of 4.
  \item \textbf{Strong joins need reads.} A join counts as strong only when at least four reads carry evidence at reliable, non-weak sites on both sides of the boundary.
  \item \textbf{Collapsed duplications and poorly mapped sites carry no joins.} A site called by more than twice the chunk's median number of reads, or by more than 1.3 times with its rarer allele under 40\% of them, or whose reads' mean mapping quality is under 40, labels reads but joins nothing.
  \item \textbf{One locus, one vote.} A local window and a call inside it, or several records of one event within 30~bp whose reads split the same way, are counted once.
  \item \textbf{No hinges.} A block that hangs on a single site with weak joins on both sides, crossed by no read, is cut at the weaker join.
\end{enumerate}

These rules leave more, shorter blocks where the reads cannot verify a join, and the reads inside each block phased correctly.

\subsection{Labels for the remaining reads, and output}

Reads the solve leaves unlabelled are labelled from evidence the solve does not vote with. Local haplotype windows (up to 2~kb) contrast the two haplotype sequences of a noisy locus, each built as the medoid of reads labelled on either side; allele windows contrast the two exact alleles of a heterozygous indel the solve does not observe, including two-allele sites with no reference read; and length windows genotype tandem repeat tracts under a stutter model learned from the chunk. A second pass of the solve uses the first two kinds as sites. A read whose calls in a block all agree takes the block's label from posterior 0.7, and a read still unlabelled is phased by its own call at one site that cannot vote, in a block of that site alone, an \emph{island} that never joins anything, so that a wrong island costs only its own reads. Chunks are joined by the reads they share, only when one orientation leads by at least three reads and 90\%; a read without a label takes one from the phase set its alleles at excluded sites associate with (one-sided binomial test, \(p \le 0.01\)). In the VCF, a phase set whose heterozygotes all lie at one locus is written unphased, as are SNPs of learned error 0.1 or more and indels of 0.25 or more.

\subsection{Evaluation}

\textbf{Data.} HG002 \citep{Zook2016} PacBio HiFi reads (DRAFT: source and chemistry) were aligned with vg Giraffe v1.73.0 \citep{Siren2021} (HiFi preset, minimizers \(k=31\), \(w=50\)) to a 16-haplotype sample \citep{Siren2024} of the Human Pangenome Reference Consortium release 2 Minigraph-Cactus graph on CHM13 \citep{Liao2023,Hickey2024,HPRC2}, from which HG002 and its parents had been removed (229 assembled samples, CHM13 and GRCh38), and surjected onto CHM13 \citep{Nurk2022}; they covered chr19 at 59-fold. The catalog was built from the whole graph. Every tool was given the same reads: pgphase their graph alignments and the surjected BAM, the other tools the BAM. HiPhase 1.6.0, LongPhase 2.0.2 and WhatsHap 2.8, by default and with \texttt{-{}-distrust-genotypes} (WhatsHap-dg), phased DeepVariant 1.10.0 calls (PACBIO model) made on the same BAM; longcallD (v0.0.11) called and phased its own. pgphase was developed on chr12, chr17, chr18, chr19 and chr20, whose truth guided its rules and constants; chr6, chr16 and chr21, and the MHC of chr6, were held out until the comparison.

\textbf{Truth.} Each read's parent of origin was found by aligning it with minimap2 \citep{Li2018} to both haplotypes of the HG002 v1.1 assembly \citep{Hansen2025} and taking the better alignment with HipHap \citep{HipHap}, which also scores the confidence of the assignment (HapQ). Truth heterozygous variants are the Q100 HG002 small-variant benchmark on CHM13 \citep{Hansen2025}.

\textbf{Read metrics.} A primary read is \emph{phaseable} when its alignment spans a heterozygous truth variant. Each phase set takes the parental orientation of the majority of its truth-labelled reads; a read is \emph{correct} when its tag agrees with it, \emph{discordant} when it does not, and \emph{unphased} when it has no tag. Reads are on the \emph{arms} or in the centromeric and pericentromeric satellites of CHM13's cenSat annotation \citep{Altemose2022} (for chr21, its acrocentric short arm too). We also report the discordant reads whose parental label is confident (HapQ $\ge$ 30): below HapQ 10, every tool is 25--37\% discordant, and most such reads carry the other parent's bases at their own clean SNVs, so their label rather than their phase is wrong.

\textbf{Variant metrics.} A tool's phased heterozygous biallelic records are matched to the truth by position, reference and alternative allele. Within a phase set, a single matched heterozygote whose orientation disagrees with its agreeing neighbours is a \emph{flip} and a lasting change a \emph{switch}. In repeats, the same haplotypes can be written as different records, and a record-level scorer then counts switches that are not. We therefore class each switch by sequence: around each of the two heterozygotes bounding it, a window of 75~bp on either side, grown until no record of either VCF crosses its edges (at most 20~kb), the tool's two haplotype sequences are compared with the truth's by edit distance, and each window pairs them in the cheaper way. A switch is \emph{real} when the two windows pair them in opposite ways, \emph{representation} when both pair them the same way, and \emph{undecided} when either window's two pairings are within 30\% (or 2 edits) of each other.
